\documentclass[12pt]{article}
\usepackage{amsmath,amssymb}
\usepackage{enumerate}
\usepackage[margin=1in]{geometry}

\begin{document}

\title{Physics 340 Final Exam: Mathematical Physics \\ Solution Set}
\author{}
\date{}
\maketitle

\section*{Question 1: Algebra of Vectors}

Given
\[
\mathbf{A}=\langle 2,-1,4\rangle,
\qquad
\mathbf{B}=\langle -1,3,2\rangle.
\]

\begin{enumerate}[a)]

\item \textbf{Dot Product}

\[
\mathbf{A}\cdot\mathbf{B}
=2(-1)+(-1)(3)+4(2)
=-2-3+8=3.
\]

\[
\boxed{\mathbf{A}\cdot\mathbf{B}=3}
\]

\item \textbf{Cross Product}

\[
\mathbf{A}\times\mathbf{B}
=
\begin{vmatrix}
\hat{i}&\hat{j}&\hat{k}\\
2&-1&4\\
-1&3&2
\end{vmatrix}
\]

\[
=
\hat{i}\bigl((-1)(2)-4(3)\bigr)
-\hat{j}\bigl(2(2)-4(-1)\bigr)
+\hat{k}\bigl(2(3)-(-1)(-1)\bigr)
\]

\[
=
-14\hat{i}-8\hat{j}+5\hat{k}.
\]

\[
\boxed{\mathbf{A}\times\mathbf{B}=\langle -14,-8,5\rangle}
\]

\item \textbf{Angle Between the Vectors}

\[
\mathbf{A}\cdot\mathbf{B}=|\mathbf{A}||\mathbf{B}|\cos\theta.
\]

\[
|\mathbf{A}|=\sqrt{21},
\qquad
|\mathbf{B}|=\sqrt{14}.
\]

\[
\cos\theta=\frac{3}{\sqrt{21}\sqrt{14}}
=\frac{3}{\sqrt{294}}.
\]

\[
\boxed{\theta=\cos^{-1}\!\left(\frac{3}{\sqrt{294}}\right)\approx79.9^\circ}
\]

\end{enumerate}

\newpage

\section*{Question 2: Coupled Linear First-Order Differential Equations}

\[
\frac{dx}{dt}=4x+y,
\qquad
\frac{dy}{dt}=2x+3y.
\]

\begin{enumerate}[a)]

\item \textbf{Matrix Form}

Let

\[
\vec{X}=
\begin{pmatrix}
x\\y
\end{pmatrix}.
\]

Then

\[
\frac{d\vec{X}}{dt}
=
\begin{pmatrix}
4&1\\
2&3
\end{pmatrix}
\vec{X}.
\]

\[
\boxed{
A=
\begin{pmatrix}
4&1\\
2&3
\end{pmatrix}}
\]

\item \textbf{Eigenvalues and Eigenvectors}

\[
\det(A-\lambda I)=0
\]

\[
\begin{vmatrix}
4-\lambda&1\\
2&3-\lambda
\end{vmatrix}
=
(4-\lambda)(3-\lambda)-2
\]

\[
=12-7\lambda+\lambda^2-2
=\lambda^2-7\lambda+10
\]

\[
(\lambda-5)(\lambda-2)=0
\]

\[
\boxed{\lambda_1=5,\qquad \lambda_2=2}
\]

For \(\lambda_1=5\):

\[
A-5I=
\begin{pmatrix}
-1&1\\
2&-2
\end{pmatrix}
\]

\[
-x+y=0 \Rightarrow y=x
\]

\[
\boxed{
\vec v_1=
\begin{pmatrix}
1\\1
\end{pmatrix}}
\]

For \(\lambda_2=2\):

\[
A-2I=
\begin{pmatrix}
2&1\\
2&1
\end{pmatrix}
\]

\[
2x+y=0 \Rightarrow y=-2x
\]

\[
\boxed{
\vec v_2=
\begin{pmatrix}
1\\-2
\end{pmatrix}}
\]

\item \textbf{Solve for \(x(t)\), \(y(t)\)}

General solution:

\[
\vec X(t)=
C_1 e^{5t}
\begin{pmatrix}
1\\1
\end{pmatrix}
+
C_2 e^{2t}
\begin{pmatrix}
1\\-2
\end{pmatrix}
\]

So

\[
x(t)=C_1e^{5t}+C_2e^{2t},
\qquad
y(t)=C_1e^{5t}-2C_2e^{2t}
\]

Use \(x(0)=1,\ y(0)=2\):

\[
C_1+C_2=1
\]

\[
C_1-2C_2=2
\]

Subtracting gives

\[
-3C_2=1
\Rightarrow
C_2=-\frac13
\]

Then

\[
C_1=\frac43
\]

Therefore

\[
\boxed{
x(t)=\frac43e^{5t}-\frac13e^{2t}}
\]

\[
\boxed{
y(t)=\frac43e^{5t}+\frac23e^{2t}}
\]

\item \textbf{Behavior of the System}

Both eigenvalues are positive, so solutions grow exponentially away from the origin. Since the eigenvalues are real and distinct, trajectories do not spiral.

The dominant long-term behavior is governed by \(e^{5t}\), so trajectories align with the eigenvector \(\langle1,1\rangle\).

\[
\boxed{\text{The origin is an unstable node.}}
\]

\end{enumerate}

\newpage

\section*{Question 3: AC RC Circuit and Complex Impedance}

\[
V(t)=V_0e^{i\omega t}
\]

A resistor \(R\) and capacitor \(C\) are connected in series.

\begin{enumerate}[a)]

\item \textbf{Total Impedance}

\[
Z_R=R,
\qquad
Z_C=\frac{1}{i\omega C}=-\frac{i}{\omega C}
\]

Thus

\[
\boxed{
Z=R-\frac{i}{\omega C}}
\]

\item \textbf{Current}

Using \(V=IZ\),

\[
I(t)=\frac{V_0e^{i\omega t}}{R-\frac{i}{\omega C}}
\]

Magnitude of impedance:

\[
|Z|=\sqrt{R^2+\left(\frac{1}{\omega C}\right)^2}
\]

So current amplitude is

\[
\boxed{
I_0=\frac{V_0}{\sqrt{R^2+\left(\frac{1}{\omega C}\right)^2}}}
\]

\item \textbf{Phase Shift}

\[
\phi_Z=
-\tan^{-1}\!\left(\frac{1}{\omega CR}\right)
\]

Therefore current has phase

\[
\boxed{
\phi=
\tan^{-1}\!\left(\frac{1}{\omega CR}\right)}
\]

Current leads voltage.

\item \textbf{Physical Meaning}

The resistor dissipates energy and produces no phase shift.  
The capacitor stores energy in the electric field, causing current to respond before voltage reaches maximum.

\end{enumerate}

\newpage

\section*{Question 4: Fourier Analysis}

\[
f(x)=
\begin{cases}
-1,&-L<x<0,\\
1,&0<x<L,
\end{cases}
\qquad \text{period }2L
\]

Use

\[
f(x)=\frac{a_0}{2}+\sum_{n=1}^{\infty}
\left[
a_n\cos\left(\frac{n\pi x}{L}\right)+
b_n\sin\left(\frac{n\pi x}{L}\right)
\right]
\]

\begin{enumerate}[a)]

\item \textbf{Coefficients}

Since \(f(x)\) is odd:

\[
\boxed{a_0=0,\qquad a_n=0}
\]

Now compute \(b_n\):

\[
b_n=\frac{1}{L}\int_{-L}^{L}f(x)\sin\left(\frac{n\pi x}{L}\right)\,dx
\]

Because odd\(\times\)odd = even,

\[
b_n=\frac{2}{L}\int_0^L \sin\left(\frac{n\pi x}{L}\right)\,dx
\]

Let

\[
k=\frac{n\pi}{L}
\]

Then

\[
b_n=\frac{2}{L}\left[-\frac{\cos(kx)}{k}\right]_0^L
\]

\[
=\frac{2}{L}\cdot\frac{1-\cos(n\pi)}{k}
\]

\[
=\frac{2}{L}\cdot\frac{1-(-1)^n}{n\pi/L}
\]

\[
\boxed{
b_n=\frac{2(1-(-1)^n)}{n\pi}}
\]

So

\[
\boxed{
b_n=
\begin{cases}
\dfrac{4}{n\pi}, & n\ \text{odd}\\[6pt]
0, & n\ \text{even}
\end{cases}}
\]

\item \textbf{Fourier Series}

\[
\boxed{
f(x)=\frac{4}{\pi}
\left[
\sin\left(\frac{\pi x}{L}\right)
+\frac13\sin\left(\frac{3\pi x}{L}\right)
+\frac15\sin\left(\frac{5\pi x}{L}\right)
+\cdots
\right]}
\]

\item \textbf{Convergence}

At points where \(f(x)\) is continuous, the Fourier series converges to \(f(x)\).

At discontinuities, it converges to the midpoint of the jump:

\[
\frac{f(x^-)+f(x^+)}{2}
\]

At \(x=0\):

\[
\frac{-1+1}{2}=0
\]

Likewise at \(x=\pm L\), the series also converges to \(0\).

\[
\boxed{\text{Continuous points: }f(x),\qquad
\text{jump points: }0}
\]

\end{enumerate}

\newpage

\section*{Question 5: Vector Calculus}

\[
\mathbf{F}=(xy)\hat{i}+(x^2+z)\hat{j}+(y+z^2)\hat{k}
\]

\begin{enumerate}[a)]

\item \textbf{Divergence}

\[
\nabla\cdot\mathbf{F}
=
\frac{\partial}{\partial x}(xy)
+\frac{\partial}{\partial y}(x^2+z)
+\frac{\partial}{\partial z}(y+z^2)
\]

\[
= y+0+2z
\]

\[
\boxed{\nabla\cdot\mathbf{F}=y+2z}
\]

\item \textbf{Curl}

\[
\nabla\times\mathbf{F}
=
\begin{vmatrix}
\hat{i}&\hat{j}&\hat{k}\\
\partial_x&\partial_y&\partial_z\\
xy&x^2+z&y+z^2
\end{vmatrix}
\]

\[
=
0\hat{i}+0\hat{j}+x\hat{k}
\]

\[
\boxed{\nabla\times\mathbf{F}=x\hat{k}}
\]

\item \textbf{Line Integral}

\[
\mathbf r(t)=\langle t,t^2,t\rangle,\qquad 0\le t\le1
\]

\[
\mathbf r'(t)=\langle1,2t,1\rangle
\]

Substitute:

\[
x=t,\ y=t^2,\ z=t
\]

\[
\mathbf F(\mathbf r(t))
=
\langle t^3,\ t^2+t,\ 2t^2\rangle
\]

Dot product:

\[
\mathbf F\cdot\mathbf r'
=
t^3+(t^2+t)(2t)+2t^2
\]

\[
=3t^3+4t^2
\]

Therefore

\[
\int_C \mathbf F\cdot d\mathbf r
=
\int_0^1(3t^3+4t^2)\,dt
\]

\[
=
\frac34+\frac43
=
\frac{25}{12}
\]

\[
\boxed{\int_C \mathbf F\cdot d\mathbf r=\frac{25}{12}}
\]

\end{enumerate}

\end{document}
